\subsection{Wuikinuxv}

Lack of FSC herring this year and
less spawn observed this year compared to the last three years.

\subsection{Kitasoo Xai'xais}

Active spawn was first observed by KX Stewardship Authority staff
in the territory at Marvin Island in Kitasu Bay on March \nth{24},
following the ponding of live fish for the commercial SOK harvest.
Natural spawn then built gradually along shorelines of
Kitasu Bay into Osment Inlet and Parson’s Anchorage, and then
along the Wilby Point shoreline for the following week.
Active spawn was also observed again this season
in Higgins Passage on March \nth{31} and
developed into a relatively extensive coverage of north and south shorelines.
Active spawn was also observed this season at several locations on
West Aristazabal Island, and
field staff were able to collect genetic samples from
Clifford Bay on April \nth{2} which is a rarely sampled location.
Several smaller spawns were observed later in the season
outside of these main spawning areas in the territory,
including several locations in Meyers Passage, as well as
a small annual deepwater spawn in Kynoch Inlet.
The Meyers Passage spawns are gradually building
following concerted effort by community harvesters
to transplant spawn on hemlock branches from larger adjacent harvesting sites.

Although sounding surveys revealed that mature herring were
fairly abundant in the area leading up to the spawning season,
community harvesters reported a relatively contracted
spawn extent, duration, and density.
Harvesters attributed this to increased humpback whale activity
in Kitasu Bay and Higgins Passage,
where up to three whales at a given time were foraging
in each of these locations.
KX Stewardship Authority staff continued to
closely monitor the piscivorous seabird community
over the herring spawn season as tens of thousands
of ducks, gulls, and alcids aggregated to forage on eggs and spawning fish.
Seabird abundance and distribution remained consistent
with observations over the past few years,
with single-day peak counts of upwards of
40,000 surf scoters observed in Kitasu Bay and Higgins Passage.
The peak surf scoter count occurred after
peak herring emergence again this season,
although thousands of gulls, scoters, and other ducks, and
hundreds of murres and loons were observed to be actively foraging
on incubating spawn across most of the region.
